\documentclass[10pt]{article}
\usepackage[margin=0.75in]{geometry}
\usepackage{amsmath,microtype}
\setlength{\emergencystretch}{3em}
\title{Guide: Turning First Contraction into Verified Descent}
\author{Collatz proof project}
\date{September 5, 2026}
\begin{document}
\maketitle
\section*{What has been proved}
For every natural starting value above one, if the powers-of-three versus
powers-of-two coefficient contracts within 65 shortcut Collatz steps,
then the actual orbit has already fallen below its starting value by that
time. This theorem covers arbitrarily large starting values.

Lean also proves that every positive return of length at most 65 reaches
one. These are finite-horizon results, not a proof of Collatz or a claim
of a new cycle-bound record.

\section*{Why coefficient contraction needs care}
After $t$ steps, with $j$ odd steps, the exact orbit value is
\[
T^t(N)=\frac{3^jN+d_t}{2^t}.
\]
Even when $3^j<2^t$, the accumulated correction $d_t$ may prevent descent.
The new argument considers the first time the coefficient contracts.
Every earlier coefficient was at least one, which bounds the correction:
$3d_t\le j3^j$.

If the endpoint still does not descend, the starting value must satisfy
\[
N\le\frac{j3^j}{3(2^t-3^j)}.
\]
For times through 65 this leaves only seeds at most 1,185. Every larger
seed descends at its first contraction in that time range.

\section*{How the remaining finite cases are checked}
A small recursive program follows each remaining seed. It succeeds as
soon as the orbit descends. Until then, it checks that the coefficient
has not contracted. Lean proves this program's soundness and checks all
seeds from two through 1,185 by ordinary kernel reduction.

The all-seed theorem combines the general bound, the checker proof, and
two kernel-checked tables. It does not assume every seed contracts or
descends within 65 steps.

\section*{What remains unresolved}
The new theorem explains the earlier-descent property of the 244
length-65 census examples for every possible seed, without certifying
the exact census list.

The correction bound is valid at any first contraction, but the finite
certificate stops at 65. We still need an argument for all later times
and for orbits with no coefficient contraction. Longer nontrivial cycles
also remain unresolved. The full Collatz goal stays open.

\section*{Verification}
The module build and 131-declaration standard-foundations audit pass.
{\raggedright Run \texttt{lake build Collatz.FirstContraction} from \texttt{lean/}.\par}
The technical paper includes the precise hypotheses and proof structure.
\end{document}
